\chapter[Probability Foundations \& Exponential Family]{Probability Foundations for Learning and the Exponential Family}
\label{chap:week11}

% ----------------------------------------------------------------------
% PART 1: EXECUTIVE ROADMAP & THE SOAP BUBBLE PARADOX
% ----------------------------------------------------------------------
\section{Executive Roadmap \& Learning Objectives}

Machine learning operates under conditions of pervasive uncertainty: observational noise, incomplete feature measurements, and stochastic data generation. Probability theory provides the formal mathematical language to quantify, model, and manipulate this uncertainty. By treating datasets as samples generated by latent joint probability distributions, we formulate estimation as statistical inference. This chapter builds continuous multivariable probability theory, analyzes non-linear transformations via Jacobian geometry, and unifies classical distributions under the powerful framework of the \textbf{Exponential Family}.

\begin{figure}[htbp]
\centering
\begin{tcolorbox}[
    colback=coralbg,
    colframe=coralborder,
    title={\large \textbf{The High-Dimensional Soap Bubble Paradox: The Vanishing Gaussian Core}},
    fonttitle=\bfseries\sffamily,
    arc=3mm,
    boxrule=1pt,
    width=\textwidth
]
\textbf{The Intuitive Fallacy:} Consider an isotropic standard Gaussian distribution in $d$ dimensions: $\mathbf{X} \sim \mathcal{N}(\mathbf{0}, I_d)$.
The probability density function is:
\begin{equation*}
f(\mathbf{x}) = \frac{1}{(2\pi)^{d/2}} \exp\left( -\frac{\|\mathbf{x}\|_2^2}{2} \right)
\end{equation*}
The mode of the distribution is at the origin $\mathbf{x} = \mathbf{0}$, where the density achieves its absolute global peak. Intuitively, one expects data samples to cluster densely around the origin, with very few points located far away in the tails.

\textbf{The Mathematical Reality:} In high dimensions ($d \ge 100$), \textbf{virtually zero probability mass resides near the origin}!
The probability of a random sample falling in a spherical shell of radius $r = \|\mathbf{x}\|_2$ is the product of the probability density and the differential volume of the hypersphere surface:
\begin{equation*}
P(r \le \|\mathbf{X}\|_2 \le r + dr) \propto r^{d-1} \exp\left( -\frac{r^2}{2} \right) dr
\end{equation*}
While the density $\exp(-r^2/2)$ decays exponentially, the surface area $S_d(r) \propto r^{d-1}$ grows explosively!

\textbf{The Soap Bubble Concentration:} The product of these two opposing forces forms a sharp spike at:
\begin{equation*}
r^* = \sqrt{d - 1} \approx \sqrt{d}
\end{equation*}
For $d = 1000$, nearly $100\%$ of all Gaussian samples lie in an ultra-thin spherical shell ("soap bubble") of radius $r \approx \sqrt{1000} \approx 31.62$ and thickness $\Delta r \approx 1$.
The center of the Gaussian is effectively a vacuum: the probability of observing a sample within distance $1$ of the mode is less than $10^{-200}$! In high dimensions, intuition forged in 2D or 3D Euclidean space fails completely.
\end{tcolorbox}
\caption{The High-Dimensional Soap Bubble Paradox: Geometry dominates density as dimension grows.}
\label{fig:soap_bubble_paradox}
\end{figure}

\subsection{Prerequisites \& Scaffolding}
Mastery of this chapter synthesizes multivariable integration, Jacobian matrices, and calculus as outlined in Table~\ref{tab:w11_prereqs}.

\begin{table}[htbp]
\centering
\small
\begin{tabularx}{\textwidth}{lXX}
\toprule
\textbf{Foundational Area} & \textbf{Core Concepts Required} & \textbf{Role in Machine Learning Foundations} \\
\midrule
Multivariable Calculus & Double/triple integrals, Jacobian determinant $|\det J|$, chain rule & Computing marginal probabilities and transforming continuous densities \\
Linear Algebra & Positive semidefinite matrices $\Sigma \succeq 0$, quadratic forms & Defining multivariate covariance matrices and Gaussian exponents \\
Elementary Probability & Random variables, Bayes' rule, independence, expectation & Formulating joint likelihoods and conditional posterior distributions \\
\bottomrule
\end{tabularx}
\caption{Prerequisite Scaffolding for Chapter 11.}
\label{tab:w11_prereqs}
\end{table}

\subsection{Core Competencies}
By the conclusion of this chapter, students will be able to:
\begin{itemize}
    \item Manipulate continuous multivariate probability density functions via marginalization and conditioning.
    \item Derive the exact transformed density $f_{\mathbf{Y}}(\mathbf{y})$ under non-linear vector mappings via the Jacobian change-of-variables theorem.
    \item Account for multiple pre-images in non-monotonic transformations (e.g., Chi-Square distributions).
    \item Formulate standard distributions (Bernoulli, Gaussian, Poisson, Exponential, Categorical) in canonical Exponential Family form.
    \item Derive expectations and covariance structures by differentiating the log-partition function $A(\boldsymbol{\eta})$.
    \item Explain the mathematical origin of the Sigmoid and Softmax functions as canonical inverse link mappings.
\end{itemize}

\subsection{Notation \& Mathematical Conventions}
\begin{itemize}
    \item $(\Omega, \mathcal{F}, P)$: A probability space with sample space $\Omega$, $\sigma$-algebra $\mathcal{F}$, and probability measure $P$.
    \item $\mathbf{X}: \Omega \to \mathbb{R}^d$: A $d$-dimensional continuous random vector.
    \item $f_{\mathbf{X}}(\mathbf{x}) \ge 0$: The joint probability density function (PDF), satisfying $\int_{\mathbb{R}^d} f_{\mathbf{X}}(\mathbf{x}) d\mathbf{x} = 1$.
    \item $f_{\mathbf{Y} \mid \mathbf{X}}(\mathbf{y} \mid \mathbf{x}) \triangleq \frac{f_{\mathbf{X}, \mathbf{Y}}(\mathbf{x}, \mathbf{y})}{f_{\mathbf{X}}(\mathbf{x})}$: The conditional density function.
    \item $J_{\mathbf{g}}(\mathbf{x}) \in \mathbb{R}^{d \times d}$: The Jacobian matrix of transformation $\mathbf{g}$.
    \item $\boldsymbol{\eta} \in \mathbb{R}^k$: Natural (canonical) parameter vector of the Exponential Family.
    \item $\mathbf{T}(x) \in \mathbb{R}^k$: Vector of sufficient statistics.
    \item $h(x) \ge 0$: Base reference measure.
    \item $A(\boldsymbol{\eta})$: The log-partition function (cumulant generating function).
\end{itemize}

% ----------------------------------------------------------------------
% PART 2: MATHEMATICAL & COMPUTATIONAL FOUNDATIONS
% ----------------------------------------------------------------------
\section{Mathematical \& Computational Foundations}

\subsection{Continuous Random Vectors and Marginalization}
Let $\mathbf{X} = [X_1, X_2, \dots, X_d]^T$ be a continuous random vector in $\mathbb{R}^d$.

\begin{definitionbox}{Joint and Marginal Densities}
The joint cumulative distribution function (CDF) is $F_{\mathbf{X}}(\mathbf{x}) = P(X_1 \le x_1, \dots, X_d \le x_d)$.
When $F_{\mathbf{X}}$ is continuously differentiable, the joint PDF is:
\begin{equation}
f_{\mathbf{X}}(\mathbf{x}) = \frac{\partial^d F_{\mathbf{X}}(x_1, \dots, x_d)}{\partial x_1 \partial x_2 \dots \partial x_d}
\end{equation}
For any partition $\mathbf{X} = (\mathbf{X}_A, \mathbf{X}_B)$ where $\mathbf{X}_A \in \mathbb{R}^{d_A}$ and $\mathbf{X}_B \in \mathbb{R}^{d_B}$, the \textbf{marginal density} is computed by integrating out the unwanted variables $\mathbf{X}_B$:
\begin{equation}
f_{\mathbf{X}_A}(\mathbf{x}_A) = \int_{\mathbb{R}^{d_B}} f_{\mathbf{X}_A, \mathbf{X}_B}(\mathbf{x}_A, \mathbf{x}_B) \, d\mathbf{x}_B
\label{eq:marginalization}
\end{equation}
\end{definitionbox}

\subsection{Transformation of Continuous Random Variables (Jacobian Rule)}

\begin{theorem}[Multivariate Change of Variables]
Let $\mathbf{X} \in \mathbb{R}^d$ have continuous PDF $f_{\mathbf{X}}(\mathbf{x})$. Let $\mathbf{Y} = \mathbf{g}(\mathbf{X})$, where $\mathbf{g}: \mathbb{R}^d \to \mathbb{R}^d$ is a bijective, continuously differentiable mapping with inverse $\mathbf{x} = \mathbf{g}^{-1}(\mathbf{y})$.
The joint PDF of $\mathbf{Y}$ is given by:
\begin{equation}
f_{\mathbf{Y}}(\mathbf{y}) = f_{\mathbf{X}}(\mathbf{g}^{-1}(\mathbf{y})) \left| \det J_{\mathbf{g}^{-1}}(\mathbf{y}) \right| = \frac{f_{\mathbf{X}}(\mathbf{g}^{-1}(\mathbf{y}))}{\left| \det J_{\mathbf{g}}(\mathbf{g}^{-1}(\mathbf{y})) \right|}
\label{eq:jacobian_transform}
\end{equation}
where $J_{\mathbf{g}}(\mathbf{x}) \in \mathbb{R}^{d \times d}$ is the Jacobian matrix with entries $[J_{\mathbf{g}}(\mathbf{x})]_{ij} = \frac{\partial g_i}{\partial x_j}$.
\end{theorem}

\begin{proof}
Let $A \subset \mathbb{R}^d$ be an arbitrary measurable Borel set in the domain of $\mathbf{X}$, and let $B = \mathbf{g}(A)$ be its image in the codomain of $\mathbf{Y}$.
By conservation of probability mass:
\begin{equation}
P(\mathbf{Y} \in B) = P(\mathbf{X} \in A) \implies \int_B f_{\mathbf{Y}}(\mathbf{y}) \, d\mathbf{y} = \int_A f_{\mathbf{X}}(\mathbf{x}) \, d\mathbf{x}
\end{equation}
By the multivariable change of variables theorem for Riemann integrals, substituting $\mathbf{x} = \mathbf{g}^{-1}(\mathbf{y})$ transforms the differential volume element:
\begin{equation}
d\mathbf{x} = \left| \det J_{\mathbf{g}^{-1}}(\mathbf{y}) \right| \, d\mathbf{y}
\end{equation}
Thus:
\begin{equation}
\int_B f_{\mathbf{Y}}(\mathbf{y}) \, d\mathbf{y} = \int_B f_{\mathbf{X}}(\mathbf{g}^{-1}(\mathbf{y})) \left| \det J_{\mathbf{g}^{-1}}(\mathbf{y}) \right| \, d\mathbf{y}
\end{equation}
Since this holds for every measurable set $B$, the integrands must agree almost everywhere, proving equation~\eqref{eq:jacobian_transform}.
\end{proof}

\begin{ponderbox}[Why the Absolute Value of the Determinant?]
Why is the absolute value $|\det J|$ mandatory in the probability density transformation rule?

\textbf{Reflection:} Probability densities are intrinsically non-negative ($f(\mathbf{y}) \ge 0$). In multivariable calculus, the determinant $\det J$ is signed: a negative determinant indicates that the transformation reverses spatial orientation (like a reflection). However, physical volume dilation is strictly positive: a reflected cube still occupies positive volume. Taking the absolute value $|\det J|$ ensures the volume scaling factor remains positive and the transformed density integrates to 1.
\end{ponderbox}

\subsection{The Exponential Family of Distributions}

\begin{definitionbox}{The Canonical Exponential Family}
A parametric family of probability distributions over random variable $X$ belongs to the \textbf{Exponential Family} if its density (or mass function) can be expressed in canonical form:
\begin{equation}
p(x; \boldsymbol{\eta}) = h(x) \exp\left( \boldsymbol{\eta}^T \mathbf{T}(x) - A(\boldsymbol{\eta}) \right)
\label{eq:exp_family_def}
\end{equation}
where:
\begin{itemize}
    \item $\boldsymbol{\eta} \in \mathbb{R}^k$: Natural (canonical) parameter vector.
    \item $\mathbf{T}(x) \in \mathbb{R}^k$: Vector of sufficient statistics.
    \item $h(x) \ge 0$: Base reference measure.
    \item $A(\boldsymbol{\eta}) \triangleq \log \int_{\mathcal{X}} h(x) \exp\left( \boldsymbol{\eta}^T \mathbf{T}(x) \right) dx$: The \textbf{log-partition function}.
\end{itemize}
\end{definitionbox}

\begin{theorem}[Moment Generating Property of the Log-Partition Function]
Let $p(x; \boldsymbol{\eta})$ belong to the Exponential Family. The gradient and Hessian of the log-partition function $A(\boldsymbol{\eta})$ generate the mean and covariance of the sufficient statistics:
\begin{align}
\nabla_{\boldsymbol{\eta}} A(\boldsymbol{\eta}) &= \E[\mathbf{T}(X)] \label{eq:exp_mean} \\
\nabla_{\boldsymbol{\eta}}^2 A(\boldsymbol{\eta}) &= \Cov(\mathbf{T}(X)) \label{eq:exp_cov}
\end{align}
Because any covariance matrix is positive semidefinite ($\Cov(\mathbf{T}(X)) \succeq 0$), the log-partition function $A(\boldsymbol{\eta})$ is strictly convex over its natural domain.
\end{theorem}

% ----------------------------------------------------------------------
% PART 3: VISUAL & GEOMETRIC INTUITION
% ----------------------------------------------------------------------
\section{Visual \& Geometric Intuition: Probability Transformations and Volume Scaling}

Figure~\ref{fig:jacobian_geometry} illustrates the geometric intuition behind the multivariate transformation rule: the non-linear map $\mathbf{g}$ deforms a rectangular differential area element $dx_1 dx_2$ into a sheared parallelogram in $\mathbf{y}$-space, scaling the local area by the absolute Jacobian determinant $|\det J_{\mathbf{g}}|$.

\begin{figure}[htbp]
\centering
\begin{tikzpicture}[
    scale=0.95,
    font=\sffamily\footnotesize
]

% Left: Domain X with orthogonal differential box
\draw[thick, fill=primarylight!40] (-4.5, 0) rectangle ++(1.8cm, 1.4cm);
\draw[->, >=stealth, color=gray] (-5.5, 0) -- (-2.0, 0) node[right] {$x_1$};
\draw[->, >=stealth, color=gray] (-4.5, -0.8) -- (-4.5, 2.2) node[above] {$x_2$};
\node[below] at (-3.6, 0) {$dx_1$};
\node[left] at (-4.5, 0.7) {$dx_2$};
\node[above] at (-3.6, 0.5) {$dA_{\mathbf{x}} = dx_1 dx_2$};
\node[below, align=center] at (-3.6, -1.2) {\textbf{Domain: } Probability Mass \\ $f_{\mathbf{X}}(\mathbf{x}) d\mathbf{x}$};

% Transformation arrow
\draw[->, >=stealth, line width=2pt, color=primarydark] (-1.5, 0.8) -- (1.0, 0.8) node[midway, above] {$\mathbf{y} = \mathbf{g}(\mathbf{x})$};

% Right: Codomain Y with deformed parallelogram
\begin{scope}[shift={(3.0, 0)}]
    \draw[->, >=stealth, color=gray] (-1.0, 0) -- (3.5, 0) node[right] {$y_1$};
    \draw[->, >=stealth, color=gray] (0, -0.8) -- (0, 2.5) node[above] {$y_2$};

    % Parallelogram spanned by Jacobian column vectors
    \filldraw[fill=tealbg!50, draw=tealborder, thick] (0, 0) -- (2.2, 0.5) -- (2.8, 2.1) -- (0.6, 1.6) -- cycle;
    \draw[->, >=stealth, line width=1.5pt, color=tealborder] (0, 0) -- (2.2, 0.5) node[midway, below] {$J_{*,1} dx_1$};
    \draw[->, >=stealth, line width=1.5pt, color=tealborder] (0, 0) -- (0.6, 1.6) node[midway, left] {$J_{*,2} dx_2$};

    \node[above] at (1.4, 0.8) {$dA_{\mathbf{y}} = |\det J_{\mathbf{g}}| dx_1 dx_2$};
    \node[below, align=center] at (1.4, -1.2) {\textbf{Codomain: } $f_{\mathbf{Y}}(\mathbf{y}) d\mathbf{y}$ \\ $= f_{\mathbf{X}}(\mathbf{x}) d\mathbf{x}$};
\end{scope}

\end{tikzpicture}
\caption{The Geometry of Probability Transformations: Conservation of probability mass requires scaling the density inversely by the local volume dilation factor $|\det J_{\mathbf{g}}|$.}
\label{fig:jacobian_geometry}
\end{figure}

% ----------------------------------------------------------------------
% PART 4: DEEP INTUITION & CONCEPTUAL DERIVATIONS
% ----------------------------------------------------------------------
\section{Deep Intuition \& Conceptual Derivations}

\subsection{Derivation of the Moment Properties of \texorpdfstring{$A(\eta)$}{A(eta)}}
By definition of a normalized probability density, the integral of $p(x; \boldsymbol{\eta})$ over the support $\mathcal{X}$ must equal 1:
\begin{equation}
\int_{\mathcal{X}} h(x) \exp\left( \boldsymbol{\eta}^T \mathbf{T}(x) - A(\boldsymbol{\eta}) \right) dx = 1 \implies \exp(A(\boldsymbol{\eta})) = \int_{\mathcal{X}} h(x) \exp\left( \boldsymbol{\eta}^T \mathbf{T}(x) \right) dx
\label{eq:exp_norm}
\end{equation}
Differentiating both sides with respect to the natural parameter $\boldsymbol{\eta}$:
\begin{equation}
\nabla_{\boldsymbol{\eta}} \exp(A(\boldsymbol{\eta})) = \exp(A(\boldsymbol{\eta})) \, \nabla_{\boldsymbol{\eta}} A(\boldsymbol{\eta}) = \int_{\mathcal{X}} h(x) \mathbf{T}(x) \exp\left( \boldsymbol{\eta}^T \mathbf{T}(x) \right) dx
\end{equation}
Dividing both sides by $\exp(A(\boldsymbol{\eta}))$:
\begin{equation}
\nabla_{\boldsymbol{\eta}} A(\boldsymbol{\eta}) = \int_{\mathcal{X}} \mathbf{T}(x) \underbrace{h(x) \exp\left( \boldsymbol{\eta}^T \mathbf{T}(x) - A(\boldsymbol{\eta}) \right)}_{p(x; \boldsymbol{\eta})} dx = \E[\mathbf{T}(X)]
\end{equation}
Differentiating a second time using the product rule:
\begin{align}
\nabla_{\boldsymbol{\eta}}^2 A(\boldsymbol{\eta}) &= \nabla_{\boldsymbol{\eta}} \left( \int_{\mathcal{X}} \mathbf{T}(x) h(x) \exp\left( \boldsymbol{\eta}^T \mathbf{T}(x) - A(\boldsymbol{\eta}) \right) dx \right) \nonumber \\
&= \int_{\mathcal{X}} \mathbf{T}(x) h(x) \exp\left( \boldsymbol{\eta}^T \mathbf{T}(x) - A(\boldsymbol{\eta}) \right) \left[ \mathbf{T}(x) - \nabla_{\boldsymbol{\eta}} A(\boldsymbol{\eta}) \right]^T dx \nonumber \\
&= \E\left[ \mathbf{T}(X) (\mathbf{T}(X) - \E[\mathbf{T}(X)])^T \right] = \Cov(\mathbf{T}(X))
\end{align}
Because the covariance matrix of any random vector satisfies $\mathbf{v}^T \Cov(\mathbf{T})\mathbf{v} = \Var(\mathbf{v}^T \mathbf{T}) \ge 0$, the Hessian $\nabla^2 A(\boldsymbol{\eta})$ is universally positive semidefinite. This proves that $A(\boldsymbol{\eta})$ is a \textbf{strictly convex functional} over its natural parameter space.

\subsection{Canonical Representation: The Bernoulli Distribution and the Sigmoid Function}
Let $X \in \{0, 1\}$ with success parameter $P(X = 1) = \mu \in (0, 1)$. The probability mass function is:
\begin{align}
p(x; \mu) &= \mu^x (1 - \mu)^{1 - x} = \exp\left( \log\left( \mu^x (1 - \mu)^{1 - x} \right) \right) \nonumber \\
&= \exp\left( x \log\mu + (1 - x)\log(1 - \mu) \right) \nonumber \\
&= \exp\left( x \log\left(\frac{\mu}{1 - \mu}\right) + \log(1 - \mu) \right)
\end{align}
Matching this expression with the canonical template $h(x) \exp\left( \eta T(x) - A(\eta) \right)$:
\begin{itemize}
    \item Base measure: $h(x) = 1$.
    \item Sufficient statistic: $T(x) = x$.
    \item Natural parameter: $\eta = \log\left(\frac{\mu}{1 - \mu}\right)$ (the \textbf{Log-Odds} or \textbf{Logit} function).
    \item Inverting for $\mu$ in terms of $\eta$:
    \begin{equation}
    e^\eta = \frac{\mu}{1 - \mu} \implies \mu = \frac{e^\eta}{1 + e^\eta} = \frac{1}{1 + e^{-\eta}} \triangleq \sigma(\eta)
    \end{equation}
    This is precisely the \textbf{logistic sigmoid activation function}!
    \item Log-partition function:
    \begin{equation}
    A(\eta) = -\log(1 - \mu) = -\log\left(1 - \frac{e^\eta}{1 + e^\eta}\right) = -\log\left(\frac{1}{1 + e^\eta}\right) = \log(1 + e^\eta)
    \end{equation}
\end{itemize}
Evaluating the derivative:
\begin{equation}
A'(\eta) = \frac{d}{d\eta} \log(1 + e^\eta) = \frac{e^\eta}{1 + e^\eta} = \frac{1}{1 + e^{-\eta}} = \mu = \E[X]
\end{equation}
Evaluating the second derivative:
\begin{equation}
A''(\eta) = \frac{d}{d\eta} \sigma(\eta) = \sigma(\eta)(1 - \sigma(\eta)) = \mu(1 - \mu) = \Var(X)
\end{equation}
This derivation reveals why the logistic sigmoid is not an arbitrary engineering heuristic: it is the mathematically unique canonical inverse link function for Bernoulli-distributed binary targets.

% ----------------------------------------------------------------------
% PART 5: MULTI-PAGE WORKED EXAMPLES
% ----------------------------------------------------------------------
\section{Multi-Page Breakable Numerical Worked Examples}

\begin{examplebox}[breakable]{Analytical Linear Transformation of a 2D Gaussian with Exact Covariance Reconstruction}
\textbf{Problem Statement:}
Let $\mathbf{X} = [X_1, X_2]^T \sim \mathcal{N}(\mathbf{0}, I_2)$ be an isotropic bivariate standard Gaussian vector with joint density:
\begin{equation}
f_{\mathbf{X}}(\mathbf{x}) = \frac{1}{2\pi} \exp\left( -\frac{1}{2}(x_1^2 + x_2^2) \right) = \frac{1}{2\pi} \exp\left( -\frac{1}{2}\mathbf{x}^T \mathbf{x} \right)
\end{equation}
Define the linear vector transformation $\mathbf{Y} = A\mathbf{X}$, where:
\begin{equation}
A = \begin{bmatrix} 2 & 1 \\ 0 & 3 \end{bmatrix}
\end{equation}

\textbf{Task Requirements:}
\begin{enumerate}
    \item Compute the determinant $\det(A)$ and the inverse matrix $A^{-1}$.
    \item Apply the Jacobian transformation theorem~\eqref{eq:jacobian_transform} to derive the exact joint PDF $f_{\mathbf{Y}}(y_1, y_2)$.
    \item Compute the theoretical covariance matrix $\Sigma_{\mathbf{Y}} = A A^T$ and its inverse $\Sigma_{\mathbf{Y}}^{-1}$.
    \item Expand the quadratic form $-\frac{1}{2}\mathbf{y}^T \Sigma_{\mathbf{Y}}^{-1}\mathbf{y}$ and verify that it matches the exponent derived via the Jacobian rule term-by-term.
\end{enumerate}

\tcblower

\textbf{Step 1: Determinant and Matrix Inversion}
The determinant is:
\begin{equation}
\det(A) = (2)(3) - (1)(0) = 6
\end{equation}
Since $\det(A) = 6 \ne 0$, the transformation is invertible.
The inverse is:
\begin{equation}
A^{-1} = \frac{1}{6} \begin{bmatrix} 3 & -1 \\ 0 & 2 \end{bmatrix} = \begin{bmatrix} \frac{1}{2} & -\frac{1}{6} \\ 0 & \frac{1}{3} \end{bmatrix}
\end{equation}

\textbf{Step 2: Change of Variables via Jacobian Transformation}
The inverse transformation mapping $\mathbf{y} \mapsto \mathbf{x}$ is:
\begin{equation}
\mathbf{x} = A^{-1}\mathbf{y} = \begin{bmatrix} \frac{1}{2}y_1 - \frac{1}{6}y_2 \\ \frac{1}{3}y_2 \end{bmatrix}
\end{equation}
The Jacobian matrix of the inverse mapping is:
\begin{equation}
J_{A^{-1}}(\mathbf{y}) = A^{-1} \implies |\det J_{A^{-1}}(\mathbf{y})| = |\det(A^{-1})| = \frac{1}{|\det(A)|} = \frac{1}{6}
\end{equation}
Now substitute $\mathbf{x} = A^{-1}\mathbf{y}$ into the original density $f_{\mathbf{X}}(\mathbf{x})$:
\begin{align}
\|\mathbf{x}\|_2^2 &= x_1^2 + x_2^2 = \left( \frac{1}{2}y_1 - \frac{1}{6}y_2 \right)^2 + \left( \frac{1}{3}y_2 \right)^2 \nonumber \\
&= \frac{1}{4}y_1^2 - \frac{2}{12}y_1 y_2 + \frac{1}{36}y_2^2 + \frac{1}{9}y_2^2 \nonumber \\
&= \frac{1}{4}y_1^2 - \frac{1}{6}y_1 y_2 + \left( \frac{1}{36} + \frac{4}{36} \right)y_2^2 \nonumber \\
&= \frac{1}{4}y_1^2 - \frac{1}{6}y_1 y_2 + \frac{5}{36}y_2^2
\end{align}
Applying the Jacobian formula:
\begin{align}
f_{\mathbf{Y}}(\mathbf{y}) &= f_{\mathbf{X}}(A^{-1}\mathbf{y}) \cdot |\det J_{A^{-1}}(\mathbf{y})| \nonumber \\
&= \frac{1}{2\pi} \exp\left( -\frac{1}{2} \left[ \frac{1}{4}y_1^2 - \frac{1}{6}y_1 y_2 + \frac{5}{36}y_2^2 \right] \right) \cdot \frac{1}{6} \nonumber \\
&= \frac{1}{12\pi} \exp\left( -\frac{1}{2} \left[ \frac{1}{4}y_1^2 - \frac{1}{6}y_1 y_2 + \frac{5}{36}y_2^2 \right] \right)
\label{eq:ex_fy_derived}
\end{align}

\textbf{Step 3: Theoretical Covariance Matrix and Its Inverse}
Since $\mathbf{X} \sim \mathcal{N}(\mathbf{0}, I_2)$, the covariance of the linear transformation $\mathbf{Y} = A\mathbf{X}$ is:
\begin{equation}
\Sigma_{\mathbf{Y}} = \E[\mathbf{Y}\mathbf{Y}^T] = \E[(A\mathbf{X})(A\mathbf{X})^T] = A \E[\mathbf{X}\mathbf{X}^T] A^T = A I_2 A^T = A A^T
\end{equation}
Performing the matrix multiplication:
\begin{equation}
\Sigma_{\mathbf{Y}} = \begin{bmatrix} 2 & 1 \\ 0 & 3 \end{bmatrix} \begin{bmatrix} 2 & 0 \\ 1 & 3 \end{bmatrix} = \begin{bmatrix} 2(2) + 1(1) & 2(0) + 1(3) \\ 0(2) + 3(1) & 0(0) + 3(3) \end{bmatrix} = \begin{bmatrix} 5 & 3 \\ 3 & 9 \end{bmatrix}
\end{equation}
Compute the determinant of $\Sigma_{\mathbf{Y}}$:
\begin{equation}
\det(\Sigma_{\mathbf{Y}}) = (5)(9) - (3)(3) = 45 - 9 = 36 = (\det A)^2
\end{equation}
The normalizing constant for a 2D Gaussian with covariance $\Sigma_{\mathbf{Y}}$ is:
\begin{equation}
\frac{1}{2\pi \sqrt{\det \Sigma_{\mathbf{Y}}}} = \frac{1}{2\pi \sqrt{36}} = \frac{1}{12\pi}
\end{equation}
This matches the normalizing prefactor in~\eqref{eq:ex_fy_derived} exactly!

Now compute the precision matrix $\Sigma_{\mathbf{Y}}^{-1}$:
\begin{equation}
\Sigma_{\mathbf{Y}}^{-1} = \frac{1}{36} \begin{bmatrix} 9 & -3 \\ -3 & 5 \end{bmatrix} = \begin{bmatrix} \frac{9}{36} & -\frac{3}{36} \\ -\frac{3}{36} & \frac{5}{36} \end{bmatrix} = \begin{bmatrix} \frac{1}{4} & -\frac{1}{12} \\ -\frac{1}{12} & \frac{5}{36} \end{bmatrix}
\end{equation}

\textbf{Step 4: Verification of the Quadratic Form}
Compute the quadratic form $\mathbf{y}^T \Sigma_{\mathbf{Y}}^{-1} \mathbf{y}$:
\begin{align}
\mathbf{y}^T \Sigma_{\mathbf{Y}}^{-1} \mathbf{y} &= \begin{bmatrix} y_1 & y_2 \end{bmatrix} \begin{bmatrix} \frac{1}{4} & -\frac{1}{12} \\ -\frac{1}{12} & \frac{5}{36} \end{bmatrix} \begin{bmatrix} y_1 \\ y_2 \end{bmatrix} \nonumber \\
&= \begin{bmatrix} y_1 & y_2 \end{bmatrix} \begin{bmatrix} \frac{1}{4}y_1 - \frac{1}{12}y_2 \\ -\frac{1}{12}y_1 + \frac{5}{36}y_2 \end{bmatrix} \nonumber \\
&= y_1\left( \frac{1}{4}y_1 - \frac{1}{12}y_2 \right) + y_2\left( -\frac{1}{12}y_1 + \frac{5}{36}y_2 \right) \nonumber \\
&= \frac{1}{4}y_1^2 - \frac{1}{12}y_1 y_2 - \frac{1}{12}y_1 y_2 + \frac{5}{36}y_2^2 \nonumber \\
&= \frac{1}{4}y_1^2 - \frac{1}{6}y_1 y_2 + \frac{5}{36}y_2^2
\end{align}
The quadratic form matches the exponent in equation~\eqref{eq:ex_fy_derived} identically. This confirms that the Jacobian transformation of a Gaussian under linear mapping $A$ produces another Gaussian $\mathcal{N}(\mathbf{0}, A A^T)$.
\end{examplebox}

\begin{examplebox}[breakable]{Non-Monotonic Transformation: Derivation of the Chi-Square Distribution}
\textbf{Problem Statement:}
Let $X \sim \mathcal{N}(0, 1)$ be a standard univariate normal random variable with PDF:
\begin{equation}
f_X(x) = \frac{1}{\sqrt{2\pi}} \exp\left( -\frac{x^2}{2} \right) \quad \text{for } x \in (-\infty, \infty)
\end{equation}
Consider the non-linear transformation:
\begin{equation}
Y = g(X) = X^2
\end{equation}

\textbf{Task Requirements:}
\begin{enumerate}
    \item Determine the support of $Y$.
    \item Identify all pre-images $x$ corresponding to a given value $y$ in the support.
    \item Compute the derivative of each inverse branch.
    \item Sum the probability contributions across all branches to derive the exact PDF $f_Y(y)$.
    \item Recognize the resulting distribution as a member of the Gamma family ($\chi^2$ with 1 degree of freedom).
\end{enumerate}

\tcblower

\textbf{Step 1: Support of $Y$}
Since $X^2 \ge 0$ for all real $X$, the support of $Y$ is strictly non-negative: $\mathcal{Y} = (0, \infty)$ (the event $Y = 0$ occurs with probability zero). For $y \le 0$, $f_Y(y) = 0$.

\textbf{Step 2: Pre-Images of the Non-Monotonic Transformation}
For any $y > 0$, the equation $x^2 = y$ has \textbf{two distinct roots}:
\begin{align}
x_1(y) &= +\sqrt{y} = y^{1/2} \\
x_2(y) &= -\sqrt{y} = -y^{1/2}
\end{align}
Because $g(x) = x^2$ is a two-to-one mapping, probability mass accumulates in $dy$ from two separate intervals in the domain of $X$: $[-\sqrt{y+dy}, -\sqrt{y}]$ and $[\sqrt{y}, \sqrt{y+dy}]$.

\textbf{Step 3: Derivatives of the Inverse Branches}
Differentiating each branch with respect to $y$:
\begin{align}
\frac{dx_1}{dy} &= \frac{d}{dy}\left( y^{1/2} \right) = \frac{1}{2}y^{-1/2} = \frac{1}{2\sqrt{y}} \\
\frac{dx_2}{dy} &= \frac{d}{dy}\left( -y^{1/2} \right) = -\frac{1}{2}y^{-1/2} = -\frac{1}{2\sqrt{y}}
\end{align}
Taking the absolute values:
\begin{equation}
\left| \frac{dx_1}{dy} \right| = \frac{1}{2\sqrt{y}}, \quad \left| \frac{dx_2}{dy} \right| = \left| -\frac{1}{2\sqrt{y}} \right| = \frac{1}{2\sqrt{y}}
\end{equation}

\textbf{Step 4: Summing Probability Contributions}
The multi-branch transformation formula is:
\begin{equation}
f_Y(y) = f_X(x_1(y))\left| \frac{dx_1}{dy} \right| + f_X(x_2(y))\left| \frac{dx_2}{dy} \right|
\end{equation}
Evaluating $f_X$ at both pre-images:
\begin{align}
f_X(+\sqrt{y}) &= \frac{1}{\sqrt{2\pi}} \exp\left( -\frac{(\sqrt{y})^2}{2} \right) = \frac{1}{\sqrt{2\pi}} e^{-y/2} \\
f_X(-\sqrt{y}) &= \frac{1}{\sqrt{2\pi}} \exp\left( -\frac{(-\sqrt{y})^2}{2} \right) = \frac{1}{\sqrt{2\pi}} e^{-y/2}
\end{align}
Substituting into the formula:
\begin{align}
f_Y(y) &= \frac{1}{\sqrt{2\pi}} e^{-y/2} \left( \frac{1}{2\sqrt{y}} \right) + \frac{1}{\sqrt{2\pi}} e^{-y/2} \left( \frac{1}{2\sqrt{y}} \right) \nonumber \\
&= 2 \cdot \frac{1}{\sqrt{2\pi}} e^{-y/2} \left( \frac{1}{2\sqrt{y}} \right) \nonumber \\
&= \frac{1}{\sqrt{2\pi}} y^{-1/2} e^{-y/2} \quad \text{for } y > 0
\end{align}

\textbf{Step 5: Connection to the Chi-Square and Gamma Distributions}
The standard Gamma PDF with shape $\alpha$ and rate $\beta$ is:
\begin{equation}
\Gamma(y; \alpha, \beta) = \frac{\beta^\alpha}{\Gamma(\alpha)} y^{\alpha - 1} e^{-\beta y}
\end{equation}
Substituting $\alpha = 1/2$ and $\beta = 1/2$, and using $\Gamma(1/2) = \sqrt{\pi}$:
\begin{equation}
\Gamma(y; 1/2, 1/2) = \frac{(1/2)^{1/2}}{\sqrt{\pi}} y^{1/2 - 1} e^{-y/2} = \frac{1}{\sqrt{2\pi}} y^{-1/2} e^{-y/2}
\end{equation}
This is the exact Chi-Square distribution with 1 degree of freedom: $Y \sim \chi^2(1)$.
\end{examplebox}

% ----------------------------------------------------------------------
% PART 6: EXPLICIT GOTCHAS & DIAGNOSTIC MISCONCEPTIONS
% ----------------------------------------------------------------------
\section{Explicit Gotchas \& Diagnostic Misconceptions}

\begin{trapbox}{Gotcha 1: Omitting the Absolute Value of the Jacobian Determinant}
A student computes a change of variables with $y = -2x$ and writes $f_Y(y) = f_X(-y/2) \cdot (-1/2)$.
\textbf{The Trap:} Probability densities are strictly non-negative: $f_Y(y) \ge 0$. The Jacobian factor represents a ratio of geometric volumes, which is always positive. The transformation formula requires the \textbf{absolute value} of the determinant: $|\det J_{\mathbf{g}^{-1}}| = |-1/2| = 1/2$. Neglecting the absolute value produces negative densities.
\end{trapbox}

\begin{trapbox}{Gotcha 2: Believing Uncorrelated Random Variables are Always Independent}
A data scientist computes sample covariance $\Cov(X, Y) = 0$ and concludes that features $X$ and $Y$ are statistically independent ($X \perp Y$).
\textbf{The Trap:} Zero covariance means there is no \textbf{linear} relationship. Non-linear dependence can be absolute. For example, let $X \sim \mathcal{U}[-1, 1]$ and $Y = X^2$. Then $\E[X] = 0$ and $\E[XY] = \E[X^3] = 0 \implies \Cov(X, Y) = 0$. However, $Y$ is completely determined by $X$; knowing $X$ completely specifies $Y$. Uncorrelatedness implies independence \textbf{if and only if} $X$ and $Y$ are jointly Gaussian.
\end{trapbox}

\begin{trapbox}{Gotcha 3: Conflating Probability Density with Probability}
A programmer evaluates a continuous Gaussian PDF and observes $f_X(0.0) = 3.989 > 1$. They declare a software bug because ``probabilities cannot exceed 1.''
\textbf{The Trap:} For continuous random variables, $f_X(x)$ is a \textbf{probability density}, not a probability. Densities have units of $1/\text{measurement}$ and can be arbitrarily large (e.g., $f(x) = 1000$ for $x \in [0, 0.001]$). The probability of an exact point is zero: $P(X = x) = 0$. Only integrals over intervals represent valid probabilities bounded in $[0, 1]$.
\end{trapbox}

\begin{trapbox}{Gotcha 4: Forgetting Multiple Pre-Images in Non-Monotonic Transformations}
An engineer computes the PDF of $Y = X^2$ for $X \sim \mathcal{N}(0, 1)$ by simply differentiating $x = \sqrt{y}$.
\textbf{The Trap:} The function $g(x) = x^2$ is not one-to-one; for each $y > 0$, there are two pre-images: $x_1 = +\sqrt{y}$ and $x_2 = -\sqrt{y}$. The transformation rule must sum over all branches:
$f_Y(y) = f_X(\sqrt{y})\left|\frac{1}{2\sqrt{y}}\right| + f_X(-\sqrt{y})\left|-\frac{1}{2\sqrt{y}}\right| = \frac{f_X(\sqrt{y})}{\sqrt{y}}$, which yields the Chi-square distribution.
\end{trapbox}

% ----------------------------------------------------------------------
% PART 7: COMPARATIVE SYNTHESIS & SUMMARY TABLE
% ----------------------------------------------------------------------
\section{Comparative Synthesis \& Summary Table}

Table~\ref{tab:exp_family_summary} summarizes the Exponential Family representations for canonical machine learning distributions.

\begin{table}[htbp]
\centering
\footnotesize
\setlength{\tabcolsep}{3pt}
\begin{tabularx}{\textwidth}{lcXX>{\raggedright\arraybackslash}Xp{2.8cm}}
\toprule
\textbf{Distribution} & \textbf{Support $\mathcal{X}$} & \textbf{Sufficient Stat $\mathbf{T}(x)$} & \textbf{Natural Param $\boldsymbol{\eta}$} & \textbf{Log-Partition $A(\boldsymbol{\eta})$} & \textbf{ML Model Link} \\
\midrule
Bernoulli($\mu$) & $\{0, 1\}$ & $x$ & $\log\left(\frac{\mu}{1 - \mu}\right)$ & $\log(1 + e^\eta)$ & Logistic Regression (Sigmoid) \\
\addlinespace
Gaussian($\mu, \sigma^2$) & $\mathbb{R}$ & $\begin{bmatrix} x \\ x^2 \end{bmatrix}$ & $\begin{bmatrix} \mu/\sigma^2 \\ -1/(2\sigma^2) \end{bmatrix}$ & $-\frac{\eta_1^2}{4\eta_2} - \frac{1}{2}\log(-2\eta_2)$ & Linear Regression (Squared Loss) \\
\addlinespace
Poisson($\lambda$) & $\mathbb{N}_0$ & $x$ & $\log\lambda$ & $e^\eta$ & Poisson Regression (Count Data) \\
\addlinespace
Exponential($\lambda$) & $[0, \infty)$ & $x$ & $-\lambda$ & $-\log(-\eta)$ & Survival Analysis, Lifetime Modeling \\
\addlinespace
Categorical($\mathbf{p}$) & $\{1, \dots, K\}$ & $[\mathbb{I}(x=k)]_{k=1}^{K-1}$ & $[\log(p_k/p_K)]_{k=1}^{K-1}$ & $\log\left(1 + \sum_{k=1}^{K-1} e^{\eta_k}\right)$ & Multi-class Softmax Regression \\
\bottomrule
\end{tabularx}
\caption{Comparative Synthesis: Canonical Exponential Family Parameterizations.}
\label{tab:exp_family_summary}
\end{table}

% ----------------------------------------------------------------------
% PART 8: PYTHON / SCIENTIFIC COMPUTING LAB
% ----------------------------------------------------------------------
\section{Python / Scientific Computing Lab}

The following production script verifies continuous transformation densities via Jacobian scaling and numerically confirms the moment-generating property of the Exponential Family log-partition function.

\begin{lstlisting}[language=Python, caption={Transformation of Densities and Exponential Family Verification in Python.}]
import numpy as np
from scipy.integrate import quad

def test_jacobian_transformation():
    """Verifies Y = exp(X) transformation: if X ~ Normal(0, 1), Y ~ LogNormal."""
    rng = np.random.RandomState(42)
    N = 100000
    x_samples = rng.randn(N)
    # Transformation: g(x) = exp(x) => g^{-1}(y) = log(y), |d/dy g^{-1}(y)| = 1/y
    y_samples = np.exp(x_samples)

    # Theoretical log-normal density: f_Y(y) = f_X(log(y)) * (1 / y)
    def theoretical_pdf(y):
        if y <= 0: return 0.0
        fx = (1.0 / np.sqrt(2.0 * np.pi)) * np.exp(-0.5 * (np.log(y)) ** 2)
        return fx * (1.0 / y)

    # Compare empirical histogram against theoretical density at y = 1.5
    y_eval = 1.5
    eps = 0.05
    empirical_density = float(np.mean((y_samples >= y_eval - eps) & (y_samples <= y_eval + eps)) / (2 * eps))
    theory_val = theoretical_pdf(y_eval)

    print("Verifying Multivariate / Jacobian Transformation:")
    print(f"  At y = {y_eval}: Empirical Density = {empirical_density:.4f}, Theoretical = {theory_val:.4f}")

def verify_exponential_family_moments():
    """Verifies that A'(eta) = E[T(X)] for the Bernoulli distribution."""
    # Bernoulli with mu = 0.75
    mu_true = 0.75
    eta = float(np.log(mu_true / (1.0 - mu_true)))

    # A(eta) = log(1 + exp(eta))
    def A(eta_val):
        return np.log(1.0 + np.exp(eta_val))

    # Numerical derivative A'(eta)
    eps = 1e-6
    grad_A = float((A(eta + eps) - A(eta - eps)) / (2.0 * eps))

    print("\nVerifying Exponential Family Moment Generation (Bernoulli):")
    print(f"  Natural Parameter eta:        {eta:.4f}")
    print(f"  True Mean E[X] = mu:          {mu_true:.4f}")
    print(f"  Derivative of Log-Partition:  {grad_A:.4f}")
    print(f"  Difference:                   {abs(grad_A - mu_true):.2e}")

def main():
    test_jacobian_transformation()
    verify_exponential_family_moments()

if __name__ == "__main__":
    main()
\end{lstlisting}

% ----------------------------------------------------------------------
% PART 9: COMPREHENSIVE PRACTICE PROBLEMS
% ----------------------------------------------------------------------
\section{Comprehensive Practice Problems}

\subsection*{Part A: Analytical \& Mathematical Problems}

\begin{enumerate}[label=\textbf{P\arabic*.}]
    \item \textbf{Analytical Linear Transformation of a 2D Gaussian:}
    Let $\mathbf{X} = [X_1, X_2]^T \sim \mathcal{N}(\mathbf{0}, I_2)$ be standard bivariate normal.
    Define linear transformation $\mathbf{Y} = A\mathbf{X}$, where $A = \begin{bmatrix} 2 & 1 \\ 0 & 3 \end{bmatrix}$.
    \begin{enumerate}[label=(\alph*)]
        \item Compute the determinant $\det(A)$ and inverse matrix $A^{-1}$.
        \item Use the Jacobian change-of-variables formula~\eqref{eq:jacobian_transform} to derive the exact joint PDF $f_{\mathbf{Y}}(y_1, y_2)$.
        \item Compute the covariance matrix $\Sigma_{\mathbf{Y}} = A A^T$ and confirm that the exponent matches $-\frac{1}{2}\mathbf{y}^T \Sigma_{\mathbf{Y}}^{-1}\mathbf{y}$.
    \end{enumerate}

    \item \textbf{Non-Monotonic Transformation (Chi-Square Distribution):}
    Let $X \sim \mathcal{N}(0, 1)$ with density $f_X(x) = \frac{1}{\sqrt{2\pi}}e^{-x^2/2}$. Let $Y = X^2$.
    \begin{enumerate}[label=(\alph*)]
        \item Identify the two pre-images $x_1(y)$ and $x_2(y)$ for $y > 0$.
        \item Compute the derivative of each inverse branch.
        \item Derive the probability density function $f_Y(y)$ and show it is a Gamma / Chi-Square distribution with 1 degree of freedom: $f_Y(y) = \frac{1}{\sqrt{2\pi}} y^{-1/2} e^{-y/2}$ for $y > 0$.
    \end{enumerate}

    \item \textbf{Exponential Distribution in Canonical Form:}
    The continuous Exponential distribution has PDF $f(x; \lambda) = \lambda e^{-\lambda x}$ for $x \ge 0$, with rate $\lambda > 0$.
    \begin{enumerate}[label=(\alph*)]
        \item Express $f(x; \lambda)$ in canonical Exponential Family form $h(x)\exp(\eta T(x) - A(\eta))$, identifying $h(x), T(x), \eta$, and $A(\eta)$.
        \item Differentiate $A(\eta)$ to compute the mean $\E[X]$.
        \item Differentiate $A(\eta)$ a second time to compute the variance $\Var(X)$.
    \end{enumerate}

    \item \textbf{Uncorrelated but Strongly Dependent Variables:}
    Let continuous random variable $X \sim \mathcal{U}[-1, 1]$. Define $Y = X^2$.
    \begin{enumerate}[label=(\alph*)]
        \item Compute $\E[X]$, $\E[Y]$, and $\E[XY]$.
        \item Show that $\Cov(X, Y) = 0$.
        \item Prove that $P(Y \le 0.25 \mid X \ge 0.8) \ne P(Y \le 0.25)$, formally proving that $X$ and $Y$ are statistically dependent despite zero covariance.
    \end{enumerate}

    \item \textbf{Derivation of the Softmax Link from Categorical Distribution:}
    Let $X \in \{1, 2, \dots, K\}$ follow a Categorical distribution with $P(X = k) = p_k$, where $\sum_{k=1}^K p_k = 1$.
    Using the reference class $K$, write the distribution in canonical Exponential Family form and prove that the inverse mapping from natural parameters $\eta_1, \dots, \eta_{K-1}$ to probabilities $p_k$ is the Softmax function:
    \begin{equation}
    p_k = \frac{e^{\eta_k}}{1 + \sum_{j=1}^{K-1} e^{\eta_j}}
    \end{equation}
\end{enumerate}

\subsection*{Part B: Diagnostic Multiple-Choice Questions}

\begin{enumerate}[label=\textbf{MCQ\arabic*.}]
    \item Let $\mathbf{Y} = A\mathbf{X} + \mathbf{b}$, where $A \in \mathbb{R}^{d \times d}$ is invertible and $\mathbf{X}$ has continuous density $f_{\mathbf{X}}(\mathbf{x})$. What is the correct probability density function $f_{\mathbf{Y}}(\mathbf{y})$?
    \begin{enumerate}[label=(\alph*)]
        \item $f_{\mathbf{Y}}(\mathbf{y}) = f_{\mathbf{X}}(A^{-1}(\mathbf{y} - \mathbf{b})) \cdot \det(A)$
        \item $f_{\mathbf{Y}}(\mathbf{y}) = f_{\mathbf{X}}(A^{-1}(\mathbf{y} - \mathbf{b})) \cdot \frac{1}{|\det(A)|}$
        \item $f_{\mathbf{Y}}(\mathbf{y}) = f_{\mathbf{X}}(A\mathbf{y} + \mathbf{b}) \cdot |\det(A)|$
        \item $f_{\mathbf{Y}}(\mathbf{y}) = f_{\mathbf{X}}(A^{-1}(\mathbf{y} - \mathbf{b}))$
    \end{enumerate}

    \item In the canonical Exponential Family $p(x; \boldsymbol{\eta}) = h(x)\exp(\boldsymbol{\eta}^T \mathbf{T}(x) - A(\boldsymbol{\eta}))$, what is the fundamental property of the Hessian $\nabla_{\boldsymbol{\eta}}^2 A(\boldsymbol{\eta})$?
    \begin{enumerate}[label=(\alph*)]
        \item It equals the negative variance of $X$ and is negative definite.
        \item It is the covariance matrix $\Cov(\mathbf{T}(X))$ and is guaranteed to be positive semidefinite, proving $A(\boldsymbol{\eta})$ is convex.
        \item It evaluates to zero for all linear models.
        \item It is equal to the Fisher information matrix multiplied by $-1$.
    \end{enumerate}
\end{enumerate}
